% ============================================================================
% ARBEITSBLATT DS 4 - Gruppe B: La vida laboral (S. 106-107)
% ============================================================================

\documentclass[11pt,a4paper]{article}

\usepackage[utf8]{inputenc}
\usepackage[T1]{fontenc}
\usepackage[ngerman]{babel}
\usepackage[left=2cm,right=2cm,top=1.8cm,bottom=1.5cm]{geometry}
\usepackage{helvet}
\usepackage[table]{xcolor}
\usepackage{tabularx}
\usepackage{array}
\usepackage{enumitem}
\usepackage{tcolorbox}
\usepackage{fancyhdr}
\usepackage{lastpage}

\renewcommand{\familydefault}{\sfdefault}

\definecolor{spanisch}{HTML}{c41e3a}
\definecolor{grau}{HTML}{666666}
\definecolor{hellgrau}{HTML}{f5f5f5}
\definecolor{gruppeB}{HTML}{1565c0}

\pagestyle{fancy}
\fancyhf{}
\fancyhead[L]{\small\textcolor{grau}{Spanisch Spätbeginner -- Gruppe B}}
\fancyhead[R]{\small\textcolor{grau}{AB-04-B}}
\fancyfoot[C]{\small\thepage/\pageref{LastPage}}
\renewcommand{\headrulewidth}{0.5pt}

\newcommand{\luecke}[1]{\underline{\hspace{#1}}}
\newcommand{\punktebox}[1]{\hfill\fbox{\small \_\_/#1}}
\newcommand{\es}[1]{\textit{#1}}

\newtcolorbox{merkkasten}[1][]{
    colback=white,
    colframe=gruppeB,
    fonttitle=\bfseries\color{gruppeB},
    title=#1,
    boxrule=1.5pt,
    arc=0pt,
    left=5pt, right=5pt, top=5pt, bottom=5pt
}

\newtcolorbox{buchverweis}{
    colback=hellgrau,
    colframe=grau,
    boxrule=0.5pt,
    arc=0pt,
    left=5pt, right=5pt, top=3pt, bottom=3pt
}

\newcounter{aufgabe}
\newcommand{\aufgabe}[2]{%
    \stepcounter{aufgabe}%
    \vspace{0.8em}%
    \noindent\textbf{\theaufgabe.} #1 \punktebox{#2}%
    \vspace{0.3em}%
}

\begin{document}

% ============================================================================
% KOPFBEREICH
% ============================================================================
\noindent
\begin{tabularx}{\textwidth}{@{}Xr@{}}
    {\Large\textbf{\textcolor{gruppeB}{La vida laboral en España}}} & Name: \luecke{5cm} \\[0.3em]
    {\small DS 4 -- Arbeitsleben + Synonyme (S. 106--107)} & Datum: \luecke{3cm} \\
\end{tabularx}

\vspace{0.3em}
\hrule
\vspace{0.8em}

% ============================================================================
% MERKKASTEN: Unterschiede DE/ES
% ============================================================================
\begin{merkkasten}[Kasten 1: Arbeitsleben -- Deutschland vs. Spanien]
\vspace{0.2em}

\begin{tabularx}{\textwidth}{@{}|X|X|@{}}
\hline
\rowcolor{hellgrau}
\textbf{Alemania} & \textbf{España} \\
\hline
horario continuo (durchgehend) & \luecke{4cm} (Mittagspause) \\
\hline
contratos indefinidos & muchos \luecke{3cm} (befristet) \\
\hline
\luecke{4cm} baja & desempleo juvenil alto \\
\hline
formación dual & menos \luecke{3cm} \\
\hline
\end{tabularx}

\vspace{0.3em}
\textbf{Konjunktionen:} \es{porque} (weil), \es{ya que} (da), \es{mientras que} (während), \es{sin embargo} (jedoch)
\vspace{0.2em}
\end{merkkasten}

\vspace{1em}

% ============================================================================
% AUFGABEN
% ============================================================================

\aufgabe{Fasse das Arbeitsleben in Spanien zusammen. (Buch S. 106, Aufg. 3b)}{8}

\begin{buchverweis}
\textbf{Schreibe 4--5 Sätze} basierend auf dem Text. Verwende \es{porque, ya que, además}.
\end{buchverweis}

\vspace{0.3em}
\begin{tabularx}{\textwidth}{@{}|X|@{}}
\hline
\\[0.4em]
\luecke{14cm} \\[0.7em]
\luecke{14cm} \\[0.7em]
\luecke{14cm} \\[0.7em]
\luecke{14cm} \\[0.7em]
\luecke{14cm} \\[0.4em]
\hline
\end{tabularx}

\vspace{0.8em}

\aufgabe{Vergleich: Erkläre 3 Unterschiede zwischen DE und ES. (Buch S. 106, Aufg. 4)}{9}

\begin{buchverweis}
\textbf{Schreibe vollständige Sätze.} Nutze: \es{mientras que, en cambio, sin embargo}.
\end{buchverweis}

\begin{enumerate}[leftmargin=2em]
    \item (Arbeitszeiten)

    \luecke{12cm}

    \luecke{12cm}

    \item (Verträge / Arbeitslosigkeit)

    \luecke{12cm}

    \luecke{12cm}

    \item (Ausbildung)

    \luecke{12cm}

    \luecke{12cm}
\end{enumerate}

\vspace{0.8em}

\aufgabe{Synonyme finden. (Buch S. 107, Aufg. 5)}{5}

\begin{buchverweis}
\textbf{Finde für jeden Begriff ein Synonym oder eine Umschreibung.}
\end{buchverweis}

\begin{tabularx}{\textwidth}{@{}|X|X|@{}}
\hline
\rowcolor{hellgrau}
\textbf{Wort} & \textbf{Synonym / Umschreibung} \\
\hline
el empleo & \luecke{5cm} \\[0.5em]
\hline
la formación profesional & \luecke{5cm} \\[0.5em]
\hline
el desempleo & \luecke{5cm} \\[0.5em]
\hline
las cualificaciones & \luecke{5cm} \\[0.5em]
\hline
la jornada laboral & \luecke{5cm} \\[0.5em]
\hline
\end{tabularx}

\vspace{0.8em}

\aufgabe{Gliederung: Vorbereitung Kurzvortrag ``Mi profesión ideal''.}{10}

\begin{buchverweis}
\textbf{Erstelle eine Gliederung} für deinen 3-4-minütigen Vortrag. Nutze Stichpunkte.
\end{buchverweis}

\vspace{0.3em}
\textbf{1. Einleitung:} Welchen Beruf? Warum interessiert dich dieser Beruf?

\luecke{12cm}

\luecke{12cm}

\vspace{0.5em}
\textbf{2. Hauptteil:} Welche Fähigkeiten brauchst du? Wie ist der Weg dorthin?

\luecke{12cm}

\luecke{12cm}

\luecke{12cm}

\vspace{0.5em}
\textbf{3. Schluss:} Warum ist dieser Beruf für dich ideal? (Persönliche Begründung)

\luecke{12cm}

\luecke{12cm}

\vspace{0.5em}
\textbf{Konjunktionen für den Vortrag:} \es{ya que, porque, además, por eso, en conclusión}

\vspace{1em}
\hrule
\vspace{0.3em}
\begin{flushright}
\textbf{Gesamt: \luecke{1cm} / 32 Punkte}
\end{flushright}

\end{document}
